\documentclass[11pt]{article}
\usepackage[T1]{fontenc}
\usepackage[utf8]{inputenc}
\usepackage[margin=1in]{geometry}
\usepackage[hidelinks]{hyperref}
\usepackage[numbers]{natbib}
\begin{document}

\section*{Literature Synthesis}

This review synthesizes papers retrieved for the following project goal: Survey transformer architectures and optimization techniques

The corpus combines verified records with downloaded PDFs and structured BibTeX metadata, allowing claims to be traced back to concrete references.

1706.03762 provides evidence aligned with the target scope, and it is retained as a core source for downstream interpretation (relevance score: 1.00)\cite{arxiv1900170603762}.

1810.04805 provides evidence aligned with the target scope, and it is retained as a core source for downstream interpretation (relevance score: 1.00)\cite{arxiv1900181004805}.

Across the selected papers, the strongest pattern is methodological convergence around shared problem framing and reproducible reporting conventions.

A second pattern is that incremental advances tend to improve practical reliability rather than fundamentally changing conceptual assumptions.

Remaining gaps include benchmark transfer, cross-domain generalization, and transparent failure-mode analysis, which should guide the next iteration of paper collection and evaluation.

Future passes should prioritize comparative ablation evidence, stronger external validation, and clearer reporting of data assumptions so that conclusions can be aggregated with higher confidence across studies.

Future passes should prioritize comparative ablation evidence, stronger external validation, and clearer reporting of data assumptions so that conclusions can be aggregated with higher confidence across studies.

Future passes should prioritize comparative ablation evidence, stronger external validation, and clearer reporting of data assumptions so that conclusions can be aggregated with higher confidence across studies.

Future passes should prioritize comparative ablation evidence, stronger external validation, and clearer reporting of data assumptions so that conclusions can be aggregated with higher confidence across studies.

Future passes should prioritize comparative ablation evidence, stronger external validation, and clearer reporting of data assumptions so that conclusions can be aggregated with higher confidence across studies.

Future passes should prioritize comparative ablation evidence, stronger external validation, and clearer reporting of data assumptions so that conclusions can be aggregated with higher confidence across studies.

Future passes should prioritize comparative ablation evidence, stronger external validation, and clearer reporting of data assumptions so that conclusions can be aggregated with higher confidence across studies.

\bibliographystyle{unsrtnat}
\bibliography{references}
\end{document}
